\documentclass[12pt]{article}


\usepackage{macros}
\usepackage{macros-gen-ro}
\usepackage{macros-curs-ro}
\usepackage{macros-math}
\usepackage{macros-logic}


\begin{document}


\renewcommand{\seminarno}{6}
\setcounter{exercitiuno}{0}

\titluseminarLM


\renewcommand{\solution}[1]{}

\parindent0pt



\parindent0pt


Pentru primele dou\u a exerci\c tii, fix\u am $\lc$ un limbaj de ordinul I care con\c{t}ine:
\bi
\item dou\u{a} simboluri de rela\c{t}ii unare $R$, $S$  \c{s}i dou\u{a} simboluri de rela\c{t}ii binare $P$, $Q$;
\item un  simbol de func\c{t}ie unar\u{a} $f$ \c{s}i un simbol de func\c{t}ie binar\u{a} $g$;
\item dou\u{a} simboluri de constante $c$, $d$.
\ei

\semex
S\u{a} se g\u{a}seasc\u{a} forme normale prenex pentru urm\u{a}toarele formule ale lui $\lc$:
\be 
\item  $\forall x(f(x)=c)\si \neg \forall z(g(y,z)=d)$;
\item  $\forall y(\forall x P(x,y)\ra \exists z Q(x,z))$;
\item $\exists x\forall y P(x,y) \sau \neg\exists y(S(y)\ra \forall z R(z))$;
\item $\exists z(\exists x Q(x, z) \sau \exists x R(x))\ra \neg (\neg \exists x R(x) \si \forall x \exists zQ(z,x))$.
\ee
\solution{
\be 
\item
\bua
\forall x(f(x)=c)\si \neg \forall z(g(y,z)=d) &\vDashv & \forall x(f(x)=c)\si \exists z \neg(g(y,z)=d)\\
&\vDashv & \forall x\exists z(f(x)=c\si \neg(g(y,z)=d)).
\eua
\item 
\bua 
 \forall y(\forall x P(x,y)\ra \exists z Q(x,z)) &\vDashv &  \forall y\exists z (\forall x P(x,y)\ra Q(x,z))\\
&\vDashv &
 \forall y\exists z (\forall u P(u,y)\ra Q(x,z))\\
&\vDashv &  \forall y\exists z \exists u (P(u,y)\ra Q(x,z)).
 \eua
\item 
 \bua
 \exists x\forall y P(x,y) \sau \neg \exists y(S(y)\ra \forall z R(z)) &\vDashv & 
  \exists x(\forall y P(x,y)\sau \neg\exists y\forall z(S(y)\ra R(z))\\
  &\vDashv &  \exists x(\forall y P(x,y)\sau \forall  y\exists z\neg (S(y)\ra R(z)) \\
&\vDashv &  \exists x(\forall u P(x,u)\sau \forall  y\exists z\neg (S(y)\ra R(z))\\
    &\vDashv &  \exists x\forall u \forall y \exists z(P(x,u)\sau \neg (S(y)\ra R(z)).
 \eua 
\item 
 \bua
\exists z(\exists x Q(x, z) \sau \exists x R(x))\ra \neg (\neg \exists x R(x) \si
\forall x \exists zQ(z,x)) &\vDashv & \\
\exists z\exists x( Q(x, z) \sau R(x))\ra (\neg\neg\exists xR(x)\sau \neg\forall x \exists zQ(z,x))
&\vDashv & \\
\exists z\exists x( Q(x, z) \sau R(x))\ra (\exists xR(x) \sau \exists x\forall z \neg Q(z,x))&\vDashv & \\
\exists z\exists x( Q(x, z) \sau R(x))\ra
\exists x (R(x) \sau \forall z \neg Q(z,x))&\vDashv & \\
\exists z\exists x( Q(x, z) \sau R(x))\ra
\exists x \forall z(R(x)\sau \neg Q(z,x))&\vDashv & \\
\exists z\exists x( Q(x, z) \sau R(x))\ra
\exists u \forall v(R(u)\sau \neg Q(v,u))&\vDashv & \\
\forall z\forall x\exists u \forall v ((Q(x, z) \sau R(x))\ra (R(u)\sau \neg Q(v,u))).
\eua
\ee
}


\semex
S\u{a} se g\u{a}seasc\u{a} o form\u a normal\u a Skolem pentru enun\c tul $\vp$ \^in form\u a 
normal\u a prenex, unde $\vp$ este, pe r\^and:
\be
\item $\forall x\exists z(f(x)=c\si \neg(g(x,z)=d))$;
\item $\forall y\exists z \exists u (P(u,y)\ra Q(y,z))$;
\item $\exists x\forall u \forall y \exists z(P(x,u)\sau \neg (S(y)\ra R(z)))$;
\item $\forall z\forall x\exists u \forall v ((Q(x, z) \sau R(x))\ra (R(u)\sau \neg Q(v,u)))$.
\ee
\solution{
\be
\item Avem $\vp^1 = \forall x(f(x)=c\si \neg(g(x,z)=d)_z(h(x)) = \forall x(f(x)=c\si \neg(g(x,h(x))=d)$, unde $h$ este un nou simbol de opera\c tie unar\u a. Cum $\vp^1$ este o formul\u a universal\u a avem $\vp^{Sk}=\vp^1$.
\item Avem $\vp^1 = \forall y \exists u (P(u,y)\ra Q(y,z))_z(p(y)) = \forall y \exists u (P(u,y)\ra Q(y,p(y)))$, unde $p$ este un nou simbol de opera\c tie unar\u a, \c si $\vp^2 = \forall y  (P(u,y)\ra Q(y,p(y)))_u(j(y)) = \forall y  (P(j(y),y)\ra Q(y,p(y)))$, unde $j$ este un nou simbol de opera\c tie unar\u a. Cum $\vp^2$ este o formul\u a universal\u a avem $\vp^{Sk}=\vp^2$.
\item Avem $\vp^1 = \forall u \forall y \exists z(P(x,u)\sau \neg (S(y)\ra R(z)))_x(m) =\forall u \forall y \exists z(P(m,u)\sau \neg (S(y)\ra R(z)))$, unde $m$ este un nou simbol de constant\u a, \c si $\vp^2 = \forall u \forall y(P(m,u)\sau \neg (S(y)\ra R(z)))_z(k(u,y)) = \forall u \forall y (P(m,u)\sau \neg (S(y)\ra R(k(u,y))))$, unde $k$ este un nou simbol de opera\c tie binar\u a. Cum $\vp^2$ este o formul\u a universal\u a avem $\vp^{Sk}=\vp^2$.
\item Avem
\begin{align*}
\vp^1 &= \forall z\forall x \forall v ((Q(x, z) \sau R(x))\ra (R(u)\sau \neg Q(v,u)))_u(n(z,x))\\
&= \forall z\forall x \forall v ((Q(x, z) \sau R(x))\ra (R(n(z,x))\sau \neg Q(v,n(z,x)))),
\end{align*}
unde $n$ este un nou simbol de opera\c tie binar\u a. Cum $\vp^1$ este o formul\u a universal\u a avem $\vp^{Sk}=\vp^1$. 
\ee 
}

\semex
Demonstra\c{t}i c\u{a} orice clas\u{a} finit axiomatizabil\u{a} $\cK$ de $\lc$-structuri este axiomatizat\u{a} de un singur enun\c{t}.

\solution{
Fie $\Gamma=\{\vp_1,\ldots, \vp_n\}$  o mul\c time finit\u a de enun\c turi  \ai $\cK=Mod(\Gamma)$. Fie $\vp:=\vp_1\wedge \vp_2 \si \ldots \si  \vp_n$. Atunci $\cA\models \Gamma \Ba \cA\models \vp_i$ pentru orice $i\in\{1,\ldots, n\}$ $\Ba$ $\cA\models \vp$. A\c{s}adar, $Mod(\vp)=Mod(\Gamma)=\cK$. 
}

\semex
Demonstra\c ti c\u a dac\u a o clas\u a $\cK$ de $\lc$-structuri  este finit axiomatizabil\u a, atunci 
complementul s\u a $\cK^c$ este de asemenea finit axiomatizabil\u a.

\solution{Conform (S6.3), $\cK$ este axiomatizat\u a de un singur enun\c t $\vp$,  
deci $\cK=Mod(\vp)$. Rezult\u a imediat c\u a pentru orice $\lc$-structur\u a $\cA$, 
\[\cA\in \cK^c \Ba \cA\notin \cK \Ba\cA\not\models \vp \Ba \cA \models \neg\vp.\] 
Prin urmare, $\cK^c=Mod(\neg\vp)$.
}

\semex
Consider\u am limbajul $\lc$ ce con\c tine un singur simbol, anume un simbol de rela\c tie de aritate $2$. S\u a se g\u aseasc\u a un enun\c t $\vp$ astfel \^inc\^at $(\Q,<) \models \vp$, dar $(\Z,<) \not\models \vp$.

\solution{
Lu\u am
$$\vp:= \forall v_0 \forall v_1 (v_0 < v_1 \to \exists v_2 (v_0 < v_2 \si v_2 < v_1)).$$

Atunci  $(\Q,<) \models \vp$ ddac\u a 
\bce 
(*) \quad pentru orice $p,q\in \Q$ cu $p<q$ exist\u a $r\in \Q$ \ai $p<r<q$,
\ece
ceea ce este adev\u arat. 

Deci, $(\Q,<) \models \vp$.

Pe de alta parte, $(\Z,<) \models \vp$ ddac\u a 
\bce 
(**) \quad pentru orice $m,n\in \Z$ cu $m<n$ exist\u a $p\in \Z$ \ai $m<p<n$,
\ece
ceea ce este fals: pentru $m=1$ \c si $n=2$, nu exist\u a  $p\in \Z$ \ai $1<p<2$. 

Deci, $(\Z,<) \not\models \vp$.
}

\semex
Fie $\lc$ un limbaj de ordinul \^int\^ai \c si $\Gamma$ o mul\c time de enun\c turi ale lui $\lc$. 
Demonstra\c ti urm\u atoarele:
\be
\item $Mod(\Gamma)=Mod(Th(\Gamma))$.
\item $Th(\Gamma)$ este o $\lc$-teorie.
\item Fie $T$ o $\lc$-teorie  \ai $\Gamma \se T$. Atunci $Th(\Gamma) \se T$.
\ee
\solution{
\be
\item "$\supseteq$"  Pentru orice $\vp\in \Gamma$, avem c\u a $\Gamma \models \vp$, deci  $\vp\in Th(\Gamma)$.  A\c sadar, $\Gamma\se Th(\Gamma)$.  Aplic\u am  Lema~4.72.(ii). \\
"$\subseteq$" Fie $\cA \in Mod(\Gamma)$ \c si $\vp\in Th(\Gamma)$ arbitrar. Avem c\u a 
$\cA \models \Gamma$ \c si $\Gamma\models \vp$, deci $\cA\models \vp$.  A\c sadar, $\cA\in  
Mod(Th(\Gamma))$.
\item   Demonstr\u am c\u a $Th(\Gamma)$ este o $\lc$-teorie. Pentru orice enun\c t $\vp$,  avem c\u a
\bua 
Th(\Gamma)\models \vp & \Ba & Mod(Th(\Gamma))\se Mod(\vp)
\Ba Mod(\Gamma)\se Mod(\vp) \text{~(conform (i))} \\
& \Ba & \Gamma\models\vp \Ba \vp \in Th(\Gamma).
\eua

\item Fie $T$ o $\lc$-teorie care con\c tine $\Gamma$ \c si $\vp\in Th(\Gamma)$. Din 
$Mod(\Gamma)\se Mod(\vp)$ \c si $Mod(T)\se Mod(\Gamma)$ rezult\u a c\u a $Mod(T)\se Mod(\vp)$, deci
$T\models \vp$. Deoarece $T$ este$\lc$-teorie, ob\c tinem c\u a $\vp\in T$. A\c sadar, $Th(\Gamma)\se T$.
\ee
}


\end{document}